\documentclass[11pt]{letter}
\usepackage[margin=1in]{geometry}
\usepackage{hyperref}

\signature{Elliot Tower}
\address{Independent researcher \\ elliot@elliottower.ai}
\date{July 7, 2026}

\begin{document}
\begin{letter}{Editors \\ \textit{PLOS ONE}}

\opening{Dear Editors,}

I am submitting ``Grassmannian geodesic distance predicts cross-cohort classifier degradation under analytical heterogeneity, after controlling for source classifier quality'' for consideration as a Research Article in \textit{PLOS ONE}.

The paper tests whether geodesic distance between cohort-specific PCA subspaces on the Grassmannian manifold predicts how much a classifier degrades when transferred to a new cohort. Raw geodesic--gap correlations are null on all seven datasets because the AUC gap conflates source classifier quality with distribution shift. After partialing out source internal AUC, geodesic distance predicts the residual gap on a colorectal cancer microbiome meta-analysis (9 studies, 824 samples; partial $\rho = +0.61$; random-forest specification CI $[+0.03, +0.77]$) and on a multi-biofluid metabolomics study (partial $\rho = +0.40$, CI $[+0.04, +0.73]$). Leave-one-study-out validation confirms out-of-sample utility (42\% MAE reduction over baseline). Five additional datasets yield null results, identifying interpretable boundary conditions: the method requires analytical heterogeneity between cohorts, and centralized platforms or shared microarray platforms produce null results even across independent study sites.

The paper owns its power limitations explicitly. A formal power analysis shows 41\% power to detect the CRC effect at the effective sample size of 9 studies; QMDiab's signal is driven by the between-type contrast (cross-biofluid vs within-biofluid) rather than graded within-type prediction. All seven datasets are public, and a pre-registration document is included. A formal Data Availability statement with repository links is provided in the manuscript.

The work meets PLOS ONE's criteria: the analyses are technically rigorous, all data and code are public, and conclusions are supported by the results. I am the sole author and declare no competing interests. The manuscript is not under consideration elsewhere.

\closing{Thank you,}

\end{letter}
\end{document}
